\section{Completion of the proof}\label{sec:completion}

\begin{proof}[Proof of Theorem~\ref{thm:main}]
In Gaussian disorder, Proposition~\ref{prop:edge-estimates} supplies the
cap estimates used in Theorem~\ref{thm:warming}. Apply that theorem to
the initial law \(\nu_n\). Proposition~\ref{prop:gaussian-subsequence}
then supplies, on a
further represented subsequence of every sequence, a normalized
positive-time density family for \(S_t^g=e^{-tc_*H_g}\), together
with locally uniform limits of both finite autocorrelations.
Proposition~\ref{prop:uniform-seed-escape} gives the averaged uniform-start
escape estimate. The sequence extraction and fixed-\(g\) uniqueness in
Theorem~\ref{thm:entrance-selection} identify every represented family
with the same measurable family \(h^g\) for almost every \(g\).
The subsequence criterion gives the Gaussian joint convergence.

The density family is bounded at each positive time, hence belongs
to \(L^2(\Pi_g)\); it is nonnegative, has mass one, and satisfies the
semigroup relation. Its weak product-topology boundary is
\(\delta_0\) by Theorem~\ref{thm:entrance-selection}.
Proposition~\ref{prop:gaussian-subsequence} proves the locally uniform
convergence and continuity of the series defining \(B_g\).
Each finite partial sum is continuous in \((s,t)\), and its evaluations
at rational pairs are measurable in \(g\) by the density and semigroup
constructions in Theorem~\ref{thm:entrance-selection}. It is therefore
Borel as a \(C_{\mathrm{loc}}((0,\infty)^2)\)-valued map, and so is its
locally uniform limit.

Proposition~\ref{prop:gaussian-subsequence} gives
\(\|h_s^g-1\|_{L^2(\Pi_g)}\to0\), and hence the required
\(L^1\) convergence. It also gives locally uniform convergence
\(B_g(s,\cdot)\to A_g(\cdot)\) for almost every \(g\) on which the
selected family is defined.
This implies the stated convergence in probability over \(g\).

Proposition~\ref{prop:universality} transfers the joint function law
to Rademacher disorder. Its comparison uses the uniform initial law
for each quench trace, so the limiting family is the one already
selected in the Gaussian argument. The added Gaussian diagonal and
the auxiliary eigenvector signs do not affect either finite
autocorrelation. For both laws the finite dynamics use rate-one site
clocks and are observed in the time unit \(T_n=n^{2/3}\). The temporary
substitution \(u=c_*s\) was reversed in the entrance proof, so the
limiting semigroup retains exactly the coefficient \(c_*\) in
\eqref{eq:edge-semigroup},
completing all the assertions.
\end{proof}
